% ==============================================================================
% SEZIONE 1: INTRODUZIONE E INQUADRAMENTO STRATEGICO
% ==============================================================================

\section{Inquadramento Strategico}

\begin{frame}{Il Nuovo Paradigma Energetico Globale}
    L'attuale mercato energetico è definito da una complessità strutturale senza precedenti che rende obsoleti i modelli di analisi tradizionali.
    \vspace{0.3cm}
    
    \note{Il mercato energetico globale affronta una complessità strutturale che rende obsoleti i modelli tradizionali.}
    
    \begin{block}{Driver di Disruption del Mercato}
        \begin{itemize}
            \item \textbf{Volatilità estrema} dei prezzi delle commodity sul mercato wholesale.
            \pause
            \note{I vecchi modelli cedono sotto tre driver di </voice><voice name="en-US-BrianNeural">disruption</voice><voice name="it-IT-DiegoNeural">. Il primo è l'estrema volatilità sul mercato </voice><voice name="en-US-BrianNeural">wholesale</voice><voice name="it-IT-DiegoNeural">.}
            
            \item \textbf{Imperativi di decarbonizzazione} e obiettivi di sostenibilità globali.
            \pause
            \note{Il secondo driver è rappresentato dagli imperativi globali di decarbonizzazione.}
            
            \item \textbf{Evoluzione tecnologica} guidata da decentralizzazione e digitalizzazione.
        \end{itemize}
    \end{block}
    \vspace{0.2cm}
    \pause
    \note{Il terzo fattore è la decentralizzazione unita alla digitalizzazione dei flussi.}
    
    \textbf{Focus Metodologico:} È necessario adottare un approccio basato su \textbf{KPI multidimensionali} che integrino metriche finanziarie (\textit{EBITDA}), parametri operativi (\textit{RAB}) e driver tecnologici.
    \note{Adottiamo un approccio basato su KPI multidimensionali, integrando la stabilità dell'EBITDA all'evoluzione della RAB.}
\end{frame}

\begin{frame}{Obiettivi del Report e Benchmark FY 2025}
    Questo documento si propone come una guida analitica, utilizzando i risultati dell'anno fiscale \textbf{FY 2025 di Enel} come benchmark industriale.
    \vspace{0.2cm}
    
    \note{Il report assume i dati di bilancio FY 2025 di Enel come benchmark di riferimento per testare il framework.}
    
    \begin{columns}[T]
        \begin{column}{0.48\textwidth}
            \begin{block}{Resilienza Finanziaria}
                Analisi della capacità di stabilizzazione dei margini e di isolamento dalle fluttuazioni esogene dei mercati.
            \end{block}
        \end{column}
        \pause
        \note{Da un lato vedremo come la stabilità dei margini dipenda dall'isolamento dal rischio merchant.}
        
        \begin{column}{0.48\textwidth}
            \begin{block}{Infrastruttura Digitale}
                La trasformazione digitale non è un mero upgrade tecnico, ma il fulcro per abilitare flussi regolati e fidelizzare l'utente.
            \end{block}
        \end{column}
    \end{columns}
    \vspace{0.3cm}
    \pause
    \note{Dall'altro, dimostreremo come la rete digitale abiliti flussi regolati, superando la logica di un semplice up grade tecnico.}
    
    \begin{alertblock}{Presupposto Fondamentale}
        La rigorosa comprensione delle variabili macroeconomiche esogene rappresenta la base di ogni valutazione sulla solidità industriale.
    \end{alertblock}
    \note{La valutazione della solidità industriale richiede una mappatura rigorosa delle variabili macroeconomiche dello scenario globale.}
\end{frame}
